\documentclass[a4paper,10pt]{article}
\usepackage[utf8]{inputenc}
\usepackage{adjustbox}
\usepackage{float}
\usepackage{indentfirst}
\usepackage{hyperref}
\usepackage{pifont}
\newcommand{\cmark}{\ding{51}}%
\newcommand{\xmark}{\ding{55}}%
\usepackage[acronym]{glossaries}

\makeglossaries
\newacronym{SPPB}{SPPB}{Simple Parallel Peripheral Bus}

%opening
\title{
	$SiER^2NA$ \\
	\small Simple Extensible Register-Register Neumann Architecture \\
	\large Simple Parallel Peripheral Bus Reference \\
	\large 1.0
}
\date{\today}

\begin{document}

\maketitle
\tableofcontents

\section{Preface}

\section{Glossary}
\printglossaries

\section{Introduction} \label{SPPB}
The \gls{SPPB} is a parallel data bus and control interface designed for connecting external devices with the CPU. It is design to be able to address a maximum of 16 connected devices at once with each connected device being identified by a 4-bit ID indicating the bus slot it is connected to.

It utilises point-to-point connections for the physical communication layer.

\section{Host Controller} \label{SPPBC}
The \acrshort{SPPB} controller is a component of the CPU that translates memory I/O requests from the MMU to electrical signals for the \acrshort{SPPB} bridge.

It internally uses a memory map that maps memory addresses to respective device ports, which is hardware-mapped as a static physical address and is accessible using standard load/store instructions.

The resulting signals from the controller are outputted via GPIO on the CPU and received by the bus bridge.

\begin{table}[H]
	\centering
	\caption{Physical interface}
	{
		\begin{tabular}{|c|c|}
\hline
GPIO pin & Description \\
\hline
0 & Data 0 \\
\hline
1 & Data 1 \\
\hline
2 & Data 2 \\
\hline
3 & Data 3 \\
\hline
4 & Data 4 \\
\hline
5 & Data 5 \\
\hline
6 & Data 6 \\
\hline
7 & Data 7 \\
\hline
8 & Clock \\
\hline
9 & Device select 0 \\
\hline
10 & Device select 1 \\
\hline
11 & Device select 2 \\
\hline
12 & Device select 3 \\
\hline
13 & Device select 4 \\
\hline
14 & Interrupt request \\
\hline
		\end{tabular}
	}
\end{table}

\subsection{Peripheral Memory Map Table} \label{SPPB:PMMT}
\begin{table}[H]
	\centering
	{
		\begin{tabular}{|c|c|}
\hline
Offset & Value \\
\hline
0x00 & Entry 0 \\
\hline
0x04 & Entry 1 \\
\hline
0x08 & Entry 2 \\
\hline
... & ... \\
\hline
0x1F8 & Entry 63 \\
\hline
\multicolumn{2}{|c|}{0x200} \\
\hline
		\end{tabular}
	}
\end{table}

\subsection{Peripheral Memory Map Entry}
\begin{table}[H]
	\centering
	\resizebox{\linewidth}{!}{
		\begin{tabular}{|c|c|c|c|c|c|c|}
\hline
63..44 & 43..24 & 23..20 & 19..16 & 15..3 & 2..1 & 0 \\
\hline
\&bs & \&lm & bid & pid & -- & rw & p \\
\hline
Physical base address[31:12] & Physical limit address[31:12] & Bus ID & Port ID & Reserved & Read/Write & Present \\
\hline
		\end{tabular}
	}
\end{table}

NOTE: Both the base and limit addresses must be \verb|0x1000| aligned.

NOTE: Base address is inclusive, while limit is exclusive, i.e. $base \le addr < limit$

NOTE: \emph{Read/write} select what I/O operation the port is used for. It can be \verb|0| for read, \verb|1| for write, or \verb|3| for both read and write.

\section{Bridge} \label{SPPBB}
The \acrshort{SPPB} bridge is, in principle, a multiplexer with a digital control protocol. It's responsible for routing memory I/O between the system board and connected peripherals.

The physical interface that connects the bridge to the system board is a 2x12-pin connector.

\begin{table}[H]
	\centering
	\caption{Physical bridge interface}
	{
		\begin{tabular}{|c|c|c|c|c|c|}
\hline
Description & \multicolumn{2}{|r|}{Pin} & \multicolumn{2}{|l|}{Pin} & Description \\
\hline
Ground & GND & 0 & 1 & D0 & Data 0 \\
\hline
Data 1 & D1 & 2 & 3 & D2 & Data 2 \\
\hline
Data 3 & D3 & 4 & 5 & GND & Ground \\
\hline
Ground & GND & 6 & 7 & D4 & Data 4 \\
\hline
Data 5 & D5 & 8 & 9 & D6 & Data 6 \\
\hline
Data 7 & D7 & 10 & 11 & GND & Ground \\
\hline
Ground & GND & 12 & 13 & CLK & Clock \\
\hline
INT select 0 & IS0 & 14 & 15 & IS1 & INT select 1 \\
\hline
INT select 2 & IS2 & 16 & 17 & IS3 & INT select 3 \\
\hline
INT select 4 & IS4 & 18 & 19 & GND & Ground \\
\hline
INT select 5 & IS5 & 20 & 21 & IS6 & INT select 6 \\
\hline
INT request & IRQ & 22 & 23 & IS7 & INT select 7 \\
\hline
		\end{tabular}
	}
\end{table}

Upon receiving a request packet from a \emph{host controller}, the bridge decodes the \emph{request frame} (See \ref{P:RF}), reads the port ID and bus ID and, using its multiplexer or similar, routes the packet to the specified device.

Afterwards, if the request packet requires a response, it waits for the device to communicate it and is then sent directly back to the \emph{host controller}.

\section{Device}
\acrshort{SPPB} devices respond to requests from the \emph{bridge} and do not send any data on their own. They are connected to the \emph{bridge} with a 2x10-pin connector.

\begin{table}[H]
	\centering
	\caption{Physical device interface}
	{
		\begin{tabular}{|c|c|c|c|c|c|}
\hline
Description & \multicolumn{2}{|r|}{Pin} & \multicolumn{2}{|l|}{Pin} & Description \\
\hline
Ground & GND & 0 & 1 & D0 & Data 0 \\
\hline
Data 1 & D1 & 2 & 3 & D2 & Data 2 \\
\hline
Data 3 & D3 & 4 & 5 & GND & Ground \\
\hline
Ground & GND & 6 & 7 & D4 & Data 4 \\
\hline
Data 5 & D5 & 8 & 9 & D6 & Data 6 \\
\hline
Data 7 & D7 & 10 & 11 & GND & Ground \\
\hline
-- & -- & 12 & 13 & -- & -- \\
\hline
Ground & GND & 14 & 15 & CLK & Clock \\
\hline
-- & -- & 16 & 17 & -- & -- \\
\hline
INT request & IRQ & 18 & 19 & GND & Ground \\
\hline
		\end{tabular}
	}
\end{table}

\section{Protocol}
\acrshort{SPPB} uses a stateless, request-response protocol. The bridge and connected devices do not send data on their own and only respond to requests by the bridge or host controller.

Every request packet is preceded by a request frame, which is a 16-bit struct that specifies basic information about the request:

\begin{table}[H] \label{P:RF}
	\centering
	\caption{Request frame structure}
	{
		\begin{tabular}{|c|c|c|c|}
\hline
15..12 & 11..8 & 7..4 & 3..0 \\
\hline
-- & pkt & pid & bid \\
\hline
Reserved & Packet ID & Port ID & Bus ID \\
\hline
		\end{tabular}
	}
\end{table}

As response packets get sent by the device and propagated to the bridge/controller only after receiving a request packet, they contain only data.
\\

NOTE: All addresses encoded into packets are relative from the mapped address base. For example, a read request to address \verb|0xA001 44FF| for a device port mapped at \verb|0xA000 0000-0xA100 0000| would be encoded and sent as \verb|0x0001 44FF|.

\subsection{Request}
\subsubsection{0x0 Read Byte Request}
\begin{table}[H]
	\centering
	{
		\begin{tabular}{|c|c|}
\hline
47..16 & 15..0 \\
\hline
adr & rf \\
\hline
Target address & Request frame (\ref{P:RF}) \\
\hline
		\end{tabular}
	}
\end{table}

\subsubsection{0x1 Read Half-Word Request}
\begin{table}[H]
	\centering
	{
		\begin{tabular}{|c|c|}
\hline
47..16 & 15..0 \\
\hline
adr & rf \\
\hline
Target address & Request frame (\ref{P:RF}) \\
\hline
		\end{tabular}
	}
\end{table}

\subsubsection{0x2 Read Word Request}
\begin{table}[H]
	\centering
	{
		\begin{tabular}{|c|c|}
\hline
47..16 & 15..0 \\
\hline
adr & rf \\
\hline
Target address & Request frame (\ref{P:RF}) \\
\hline
		\end{tabular}
	}
\end{table}

\subsubsection{0x3 Write Byte}
This request does not have a response packet

\begin{table}[H]
	\centering
	{
		\begin{tabular}{|c|c|c|}
\hline
55..48 & 47..16 & 15..0 \\
\hline
d & adr & rf \\
\hline
Data & Target address & Request frame (\ref{P:RF}) \\
\hline
		\end{tabular}
	}
\end{table}

\subsubsection{0x4 Write Half-Word}
This request does not have a response packet

\begin{table}[H]
	\centering
	{
		\begin{tabular}{|c|c|c|}
\hline
63..48 & 47..16 & 15..0 \\
\hline
d & adr & rf \\
\hline
Data & Target address & Request frame (\ref{P:RF}) \\
\hline
		\end{tabular}
	}
\end{table}

\subsubsection{0x5 Write Word}
This request does not have a response packet

\begin{table}[H]
	\centering
	{
		\begin{tabular}{|c|c|c|}
\hline
79..48 & 47..16 & 15..0 \\
\hline
d & adr & rf \\
\hline
Data & Target address & Request frame (\ref{P:RF}) \\
\hline
		\end{tabular}
	}
\end{table}

\subsubsection{0x6 Interrupt Information Request} \label{P:IIRQ}
\begin{table}[H]
	\centering
	{
		\begin{tabular}{|c|c|c|}
\hline
15..0 \\
\hline
rf \\
\hline
Request frame (\ref{P:RF}) \\
\hline
		\end{tabular}
	}
\end{table}

\subsection{Response}
\subsubsection{0x0 Read Byte Response}
\begin{table}[H]
	\centering
	{
		\begin{tabular}{|c|}
\hline
7..0 \\
\hline
d \\
\hline
Data \\
\hline
		\end{tabular}
	}
\end{table}

\subsubsection{0x1 Read Half-Word Response}
\begin{table}[H]
	\centering
	{
		\begin{tabular}{|c|}
\hline
15..0 \\
\hline
d \\
\hline
Data \\
\hline
		\end{tabular}
	}
\end{table}

\subsubsection{0x2 Read Word Response}
\begin{table}[H]
	\centering
	{
		\begin{tabular}{|c|}
\hline
31..0 \\
\hline
d \\
\hline
Data \\
\hline
		\end{tabular}
	}
\end{table}

\subsubsection{0x6 Interrupt Information Response}
\begin{table}[H]
	\centering
	{
		\begin{tabular}{|c|c|c|}
\hline
23..16 & 15..0 \\
\hline
int & rf \\
\hline
Interrupt ID & Request frame (\ref{P:RF}) \\
\hline
		\end{tabular}
	}
\end{table}

If the device did not raise any interrupt, \verb|int| is 0.

\section{Interrupts}
Interrupts are first handled within the device-bridge loop before being sent to the CPU as a binary selector.

Whenever a connected device wants to raise an interrupt, it communicates the need to do so by pulling the \verb|IRQ| pin HIGH on the bridge. After every operation the \emph{bridge} checks the status of the \verb|IRQ| bin and, if it's detected HIGH, it sents out a \hyperref[P:IIRQ]{Interrupt Information Request} packet to every connected device in an unspecified order, propagating the interrupt request to the CPU after decoding the response packet and if \verb|int| is not zero. The \verb|IRQ| pin is then cleared by the requesting device after receiving the IIRQ packet.

% TODO there are kinda 2 ways to do this; this one would mean that if some device is constantly raising interrupts then it will stall/drown the others unless the bridge uses some kind of random queue instead of going from 0 to 15, the other option is going through the loop until IRQ is cleared, but that means that a rogue device could stall the bridge and system forever. all-in-all this is probably the better option and has a simple way to solve it
If \verb|IRQ| is still HIGH after receiving a non-zero interrupt ID, i.e. more than one device wants to raise an interrupt at the same time, the \emph{bridge} completes a single request from the \emph{host controller} before going through the same process.

\end{document}
